\section{Residual correction: identities, comparisons and limits}\label{sec:identity}

The learning curves motivate a distinction between an identity for a fitted predictor and a
claim about its out-of-sample rate. The former can be established without distributional
assumptions; the latter requires assumptions not measured by these experiments. We first fix
all fitted quantities, including any data-dependent kernel hyperparameters, and work with one
output coordinate. The identities extend coordinatewise to vector outputs.

\subsection{A frozen regression operator}
Fix a design $X_n=(x_1,\ldots,x_n)$, a positive-semidefinite kernel with Gram matrix $K$, and
$\lambda>0$. Define
\begin{equation}\label{eq:krr}
 (Sg)(x)=k(x,X_n)(K+\lambda I)^{-1}g(X_n),\qquad R=I-S.
\end{equation}
Here $\lambda$ is the diagonal regularizer in the linear system. The campaign code uses
$\lambda=n\gamma$, where $\gamma$ is its reported nugget; the distinction matters when the
training size changes. Although a selected kernel can depend on the targets, $S$ is linear
in the response only \emph{after that selection is frozen}.

\begin{proposition}[Frozen residual correction]\label{prop:identity}
Let $f$ be the target and $h$ any fitted mean, and set
$\widehat f_c=h+S(f-h)$. Then:
\begin{align}
 f-\widehat f_c&=R(f-h),\label{eq:residual-exact}\\
 (Rg)(X_n)&=\lambda(K+\lambda I)^{-1}g(X_n).\label{eq:design-residual}
\end{align}
If $K u_j=\mu_j u_j$, the design residual in direction $u_j$ is multiplied by
$\lambda/(\mu_j+\lambda)\in(0,1]$. The factor is one when $\mu_j=0$ and strictly below one
when $\mu_j>0$; this is shrinkage, not a hard spectral projection.

Suppose, in addition, that $\mathcal S$ is a normed vector space containing $f$ and $h$ and
that the following bound holds \emph{uniformly} on that space for the realized operator:
\[
 \|Rg\|_{L^2(P)}\le C_n\|g\|_{\mathcal S}\quad(g\in\mathcal S).
\]
Then
\begin{equation}\label{eq:inherit}
 \|f-\widehat f_c\|_{L^2(P)}\le C_n\|f-h\|_{\mathcal S}.
\end{equation}
Consequently, $\|f-h\|_{\mathcal S}<\|f\|_{\mathcal S}$ makes this displayed upper bound
smaller than $C_n\|f\|_{\mathcal S}$ when $C_n>0$. It is neither a necessary nor a
sufficient condition for the corrected predictor's actual error to be smaller than that of
$Sf$.
\end{proposition}
\begin{proof}
Subtract $h+S(f-h)$ from $f$ to obtain \eqref{eq:residual-exact}. On the design,
$I-K(K+\lambda I)^{-1}=\lambda(K+\lambda I)^{-1}$, which proves
\eqref{eq:design-residual} and its spectral form. Apply the assumed uniform inequality to
$g=f-h$ for \eqref{eq:inherit}. The final distinction is demonstrated below.
\end{proof}

\paragraph{Why the bounds do not order the errors.}
On a two-point design take $K=\operatorname{diag}(1,99)$ and $\lambda=1$, so
$R=\operatorname{diag}(1/2,1/100)$ and $C_n=1/2$ is its Euclidean operator norm. With
$f=(1,0)$ and $f-h=(0,2)$, the residual norm is twice the target norm, yet
$\|R(f-h)\|=0.02<0.50=\|Rf\|$. Conversely, with $f=(0,1)$ and $f-h=(1/2,0)$, the
residual norm is smaller, but $\|R(f-h)\|=0.25>0.01=\|Rf\|$. These examples also show why
a more accurate neural mean need not yield a more accurate corrected predictor: the
orientation of the remaining error matters.

\subsection{An exact criterion for improvement}
The relevant comparison can be made exactly in any Hilbert error metric. This includes a
finite weighted squared error and $L^2(P)$, but not directly the mean of per-sample relative
norms or the nonlinear radiance and reflectance metrics in Table~\ref{tab:results}.

\begin{proposition}[Error alignment]\label{prop:alignment}
For the same frozen operator $S$, put $a=Rf$ and $b=Rh$. Then
\begin{equation}\label{eq:alignment}
 \|f-\widehat f_c\|^2-\|f-Sf\|^2=\|b\|^2-2\langle a,b\rangle.
\end{equation}
Thus full correction improves on $Sf$ exactly when
$2\langle a,b\rangle>\|b\|^2$. More generally, the interpolation
$p_t=Sf+tRh$, $0\le t\le1$, has minimum squared error at
\begin{equation}\label{eq:shrink-mean}
 t_* = \min\{1,\max\{0,\langle a,b\rangle/\|b\|^2\}\}
 \quad\text{if }b\ne0;
\end{equation}
if $b=0$, every $t$ is equivalent. If the direct fit uses $S_0$ and the correction uses
$S_1$, \eqref{eq:alignment} instead holds with
$a=(I-S_0)f$ and $b=(S_1-S_0)f+(I-S_1)h$.
\end{proposition}
\begin{proof}
For a common operator the corrected error is $a-b$, so expansion of its squared norm gives
\eqref{eq:alignment}. The error of $p_t$ is $a-tb$; minimize the resulting quadratic on
$[0,1]$ to obtain \eqref{eq:shrink-mean}. For distinct operators,
$f-h-S_1(f-h)=(I-S_0)f-[(S_1-S_0)f+(I-S_1)h]$, after which the same expansion applies.
\end{proof}

The criterion is a diagnostic, not an oracle available at prediction time. Replacing its inner
products by validation errors gives a validation-optimal interpolation, not a guarantee on new
data. The reported experiments use separately tuned direct and residual kernels, so their
comparison must use the final, two-operator form rather than silently identifying the operators.
Neither the criterion nor the training-design contraction establishes that correction will
beat both parents at every sample size.

\subsection{What the principal-component implementation computes}
The campaign trains a network on all standardized bands but projects its predictions into
the retained principal-component basis before correction. Thus its corrected predictor is
not, literally, the unprojected network plus a kernel residual. To state the implementation
precisely, let $f$ and $h$ now denote the target and network output after the common fitted
standardization and PCA centering. Let $\Pi$ be the orthogonal projection onto the retained
spectral coordinates. In these centered coordinates, the implemented prediction is
\[
 \widehat f_{c,\Pi}=\Pi h+S\Pi(f-h).
\]
The output mean and diagonal band scaling are restored afterwards. Since the same scalar
kernel operator acts on each retained output coordinate,
\begin{equation}\label{eq:projected-residual}
 f-\widehat f_{c,\Pi}=(I-\Pi)f+(I-S)\Pi(f-h).
\end{equation}
The terms on the right are orthogonal in the standardized spectral Euclidean norm at each
query, hence their squared norms add there. They need not be orthogonal in physical units
after unequal band rescaling. This decomposition separates truncation error from residual
regression and explains why output rank is a genuine part of the model specification. It
does not by itself prove the observed optimum near rank 32. Comparisons with the unprojected
network include the effect of projection as well as residual fitting.

\paragraph{Rates and interpretation.}
Standard kernel learning-rate results impose source, spectrum, sampling and regularization
conditions \cite{cdv,sutherland,fs}. If a bound uniform over the relevant residual class holds
with $C_n=O(n^{-\alpha})$ and one \emph{separately establishes}
$\|f-h_n\|_{\mathcal S}=O(n^{-\beta})$, then \eqref{eq:inherit} gives an upper bound
$O(n^{-(\alpha+\beta)})$. It gives neither equality of rates nor a lower bound. Empirical
network loss does not establish convergence in $\mathcal S$, and a bound for a fixed target
cannot automatically be applied to an adaptively fitted residual or a selected kernel. The
required smoothness norms and uniform events are not measured here. Accordingly the slopes
in Table~\ref{tab:slopes} are finite-range empirical slopes, not estimates of a proved
asymptotic exponent. Smoother or better-aligned residuals are possible interpretations of the
successful corrections; they are not conclusions forced by the algebra.

The tests in \path{code/test_publication_revision.py} check the residual identities,
regularizer convention, semidefinite boundary case, alignment formula, counterexamples and
projected implementation on explicit and randomized finite designs. These numerical checks
support the implementation of the displayed formulas; the proofs above supply their
mathematical justification.
